\section{テフラ同定・年代対比と再堆積}

\subsection{平成28年度 問34}
\label{q:h28-reg-applied-34}

\noindent\textbf{原典：} sources/H28/17\_applied\_science.pdf、PDF物理13頁。

\textbf{公式正答：} 3（single）。

\par\textbf{検討未確定：} 海成層が常に鍵層、火山灰が常に一番確かという強い一般化は成立範囲を要する。花粉の淡水限定の誤りは確認済みだが、他肢の代表性を部分校正で追跡する。

\paragraph{問題の要旨} 地層対比で使う微化石・花粉・火山灰・古地磁気の情報と適用域を比較する。

\Needspace{4\baselineskip}\paragraph{確認した条件・記述（独自要約）}\mbox{}\par

\noindent\textbf{判断方向：} 最も不適切・誤りを選ぶ\par

\noindent\textbf{原典照合：} 15原典画像で18問の問題幹、全5肢、否定、専門語、化学式、候補順を確認。\par

\noindent\textbf{適用条件：} 流れと構造の典型対応は一意成因を保証しない。用語、時間間隙、保存、試料処理の条件を個別に示す。\par

\paragraph{問われる核心} 地層対比で使う微化石・花粉・火山灰・古地磁気の情報と適用域を比較する。

\paragraph{必要知識} 花粉の運搬・海成保存、微化石群集、鍵層と時間面、初生テフラ、定方位試料。

\paragraph{考え方・正答を選ぶ理由（査読待ち初稿）} 花粉は陸上で作られても風や河川によって海へ運ばれる。海成層の花粉分析による層序対比が実際に報告されているため、淡水成層だけとする3が最も不適切である。微化石と古地磁気は別の観測量を持ち、火山灰は一回の噴火を同定できる場合の対比に強い。広い海成層や火山灰の利点を、全条件での同時性・最確性へ広げない。

\Needspace{4\baselineskip}\paragraph{選択肢の検討（査読待ち初稿）}\mbox{}\par

\noindent\textbf{1：適切（条件を含めて検討）} 海成層の広い分布を鍵層として使う。広く追跡できる特徴的な海成層は鍵となることがある。ただし海成であるだけで同時面とはならない。\par

\noindent\textbf{2：適切（条件を含めて検討）} 微化石を堆積環境調査に使う。珪藻、有孔虫、放散虫、貝形虫などの組成が生息域や堆積環境の材料となる。再堆積も考慮する。\par

\noindent\textbf{3：誤り・最も不適切} 花粉は淡水成層のみに含まれる。海成層でも花粉化石が記録され、陸上植生の変遷や対比に使われる。「のみ」という限定が誤りである。\par

\noindent\textbf{4：適切（条件を含めて検討）} 同一火山灰層を最も確かな対比の決め手とする。同一噴火の初生テフラを正しく同定すれば強い時間対比となる。一律にどの手法より確かとはせず、再堆積と誤対比を除く。\par

\noindent\textbf{5：適切（条件を含めて検討）} 古地磁気は定方位試料で調べる。磁化方向を地理方向へ対応させる必要があるため方位を保持する。傾動や再磁化の補正も別に必要である。\par

\paragraph{再利用する考え方} 観察できる構造と成因・環境の解釈を区別する。代表モデルに適用条件を添え、単一の形や測定値から唯一の環境を断定しない。

\paragraph{共通理論} \hyperref[theory:stratigraphy:marker]{第\ref*{theory:stratigraphy:marker}節}, \hyperref[theory:stratigraphy:pollen]{第\ref*{theory:stratigraphy:pollen}節}, \hyperref[theory:stratigraphy:scope]{第\ref*{theory:stratigraphy:scope}節}

\par 参考：\cite{strat_pollen_tokushima,strat_gsj_atsumi}。

\paragraph{関連問題} \hyperref[q:r02-reg-applied-31]{R02-REG-APPLIED-31}

\par 難易度 3（校準初稿）、優先度未評価。

\subsection{令和2年度 問31}
\label{q:r02-reg-applied-31}

\noindent\textbf{原典：} sources/R02/17\_applied\_science.pdf、PDF物理11頁。

\textbf{公式正答：} 4（single）。

\paragraph{問題の要旨} 火山砕屑物を使う地層対比の適用範囲と、野外記載・鉱物分析の役割を比較する。

\Needspace{4\baselineskip}\paragraph{確認した条件・記述（独自要約）}\mbox{}\par

\noindent\textbf{判断方向：} 最も適切な記述を選ぶ\par

\noindent\textbf{原典照合：} 15原典画像で18問の問題幹、全5肢、否定、専門語、化学式、候補順を確認。\par

\noindent\textbf{適用条件：} 流れと構造の典型対応は一意成因を保証しない。用語、時間間隙、保存、試料処理の条件を個別に示す。\par

\paragraph{問われる核心} 火山砕屑物を使う地層対比の適用範囲と、野外記載・鉱物分析の役割を比較する。

\paragraph{必要知識} 破片状火山物質、局地/広域の用途、海湖成保存、中期更新世対比例、野外・分析の併用。

\paragraph{考え方・正答を選ぶ理由（査読待ち初稿）} テフラは火山の破片状噴出物を扱うので溶岩と区別する4が適切である。他は有効な材料を不当に除外する。局地層も局地対比に使え、海成・湖成層も対象となり、中期更新世の対比例がある。同定は野外の形質とガラス/鉱物分析を併用する。テフラの狭義・広義が文献で違うこと、初生と再堆積が年代解釈を変えることは別の確認事項である。

\Needspace{4\baselineskip}\paragraph{選択肢の検討（査読待ち初稿）}\mbox{}\par

\noindent\textbf{1：非該当・不適切} 局地的テフラを避け広域テフラだけを使う。局地的噴火層もその地域の時間対比に使える。広域性は用途の利点であり必須条件ではない。\par

\noindent\textbf{2：非該当・不適切} 陸上地層だけに適用する。海や湖の堆積物中のテフラも時間対比に使える。\par

\noindent\textbf{3：非該当・不適切} 色・粒度・層厚を使わず屈折率と組成で同定する。分析値は重要だが野外の層序位置・色・粒径・層厚も使う。GSJの対比訂正には粒径/層厚による候補排除がある。\par

\noindent\textbf{4：該当・適切} テフラは一般に火山砕屑物で溶岩と区別する。破片状の噴出物を連続した溶岩流と分けるという本問の一般的な用語に対応する。\par

\noindent\textbf{5：非該当・不適切} 後期更新世以降だけに適用する。GSJの中期更新世海成層の対比例だけでもこの限定への反例になる。古い層では保存と変質を別途検討する。\par

\paragraph{再利用する考え方} 観察できる構造と成因・環境の解釈を区別する。代表モデルに適用条件を添え、単一の形や測定値から唯一の環境を断定しない。

\paragraph{共通理論} \hyperref[theory:stratigraphy:tephra]{第\ref*{theory:stratigraphy:tephra}節}, \hyperref[theory:stratigraphy:identify]{第\ref*{theory:stratigraphy:identify}節}, \hyperref[theory:stratigraphy:scope]{第\ref*{theory:stratigraphy:scope}節}

\par 参考：\cite{volc_usgs_tephra,strat_gsj_tephra2013,strat_gsj_atsumi}。

\paragraph{関連問題} \hyperref[q:h28-reg-applied-34]{H28-REG-APPLIED-34}

\par 難易度 3（校準初稿）、優先度未評価。

\section{不整合の時間間隙と三分類}

\subsection{令和7年度 問28}
\label{q:r07-reg-applied-28}

\noindent\textbf{原典：} sources/R07/17\_applied\_science.pdf、PDF物理10頁。

\textbf{公式正答：} 1（single）。

\paragraph{問題の要旨} 地層境界の時間間隙と、基盤岩種・層理の幾何による不整合の分類を比較する。

\Needspace{4\baselineskip}\paragraph{確認した条件・記述（独自要約）}\mbox{}\par

\noindent\textbf{判断方向：} 最も不適切・誤りを選ぶ\par

\noindent\textbf{原典照合：} 15原典画像で18問の問題幹、全5肢、否定、専門語、化学式、候補順を確認。\par

\noindent\textbf{適用条件：} 流れと構造の典型対応は一意成因を保証しない。用語、時間間隙、保存、試料処理の条件を個別に示す。\par

\paragraph{問われる核心} 地層境界の時間間隙と、基盤岩種・層理の幾何による不整合の分類を比較する。

\paragraph{必要知識} 層序の時間間隙、侵食と非堆積、傾動、基盤岩種、層理の走向・傾斜。

\paragraph{考え方・正答を選ぶ理由（査読待ち初稿）} 不整合面を挟む間に記録されない時間と侵食がある。下位層の堆積後、固結・傾動・侵食が起きてから上位層が堆積することもあるため、下位の堆積直後の地形とする1が最も不適切である。無整合は基盤の火成・変成岩との関係、平行と傾斜不整合は上下層理の幾何で分ける。平行であるだけで整合とする判断を避ける。

\Needspace{4\baselineskip}\paragraph{選択肢の検討（査読待ち初稿）}\mbox{}\par

\noindent\textbf{1：誤り・最も不適切} 不整合面は多くの場合下位層の堆積直後の地形面。下位堆積後の非堆積、傾動や侵食を経て上位がたまるため、下位の堆積直後の面とは限らない。\par

\noindent\textbf{2：適切（条件を含めて検討）} 層状岩が非層状火成・変成岩を覆う不整合を無整合とする。堆積層が火成・変成基盤を不整合に覆う関係に対応する。\par

\noindent\textbf{3：適切（条件を含めて検討）} 平行な上下層で下位が顕著に削剥される境界を平行不整合とする。層理の平行性は時間的連続を保証しない。侵食と時間間隙がある平行不整合である。\par

\noindent\textbf{4：適切（条件を含めて検討）} 上下層の走向・傾斜が違う不整合を傾斜不整合とする。下位が傾動・侵食された後に上位が堆積する幾何に対応する。\par

\noindent\textbf{5：適切（条件を含めて検討）} 境界に時間間隙がほぼない関係を整合とする。本問の層序関係の定義である。検出できない短い休止までゼロと断定する意味ではない。\par

\paragraph{再利用する考え方} 観察できる構造と成因・環境の解釈を区別する。代表モデルに適用条件を添え、単一の形や測定値から唯一の環境を断定しない。

\paragraph{共通理論} \hyperref[theory:stratigraphy:hiatus]{第\ref*{theory:stratigraphy:hiatus}節}, \hyperref[theory:stratigraphy:unconformity]{第\ref*{theory:stratigraphy:unconformity}節}

\par 参考：\cite{strat_iu_unconformity}。

\par 難易度 3（校準初稿）、優先度未評価。
